% =============================================================
% DOCUMENTCLASS
% =============================================================
\documentclass[12pt,oneside]{book}

% =============================================================
% METADATOS
% =============================================================
\newcommand{\universidad}{Universidad de Buenos Aires}
\newcommand{\materia}{Derecho Procesal Civil}
\newcommand{\titulo}{Sistemas de Solución de Conflictos: Métodos Alternativos de Resolución de Conflictos}
\newcommand{\autor}{Isaac Edgar Camacho Ocampo}
\newcommand{\anio}{2026}

% =============================================================
% PAQUETES
% =============================================================
\usepackage[utf8]{inputenc}
\usepackage[T1]{fontenc}
\usepackage[spanish,es-nodecimaldot]{babel}

\usepackage[
    top=2.5cm,
    bottom=2.5cm,
    left=3cm,
    right=2.5cm,
    headheight=15pt
]{geometry}

\usepackage{amsmath}
\usepackage{amssymb}
\usepackage{mathtools}

\usepackage{graphicx}
\usepackage{float}
\usepackage{caption}
\usepackage{subcaption}

\usepackage{booktabs}
\usepackage{array}
\usepackage{longtable}
\usepackage{tabularx}

\usepackage{enumitem}

\usepackage{xcolor}
\usepackage[most]{tcolorbox}

\usepackage{fancyhdr}

\usepackage{ragged2e}

% =============================================================
% RUTAS DE IMÁGENES
% =============================================================
\graphicspath{
    {./img/}
    {./graficos/}
    {./figuras/}
}

% =============================================================
% HIPERVÍNCULOS Y METADATOS PDF
% =============================================================
\usepackage[
    colorlinks=true,
    linkcolor=blue!50!black,
    citecolor=green!40!black,
    urlcolor=blue!60!black,
    pdftitle={\titulo},
    pdfauthor={\autor},
    pdfsubject={Apuntes universitarios},
    bookmarks=true
]{hyperref}

% =============================================================
% CONFIGURACIÓN
% =============================================================
\setcounter{tocdepth}{2}

\setlength{\parindent}{0pt}
\setlength{\parskip}{0.5em}

\setlist[itemize]{
    topsep=0.3em,
    itemsep=0.2em
}

\setlist[enumerate]{
    topsep=0.3em,
    itemsep=0.2em
}

\clubpenalty=10000
\widowpenalty=10000

% =============================================================
% ENCABEZADOS Y PIES
% =============================================================
\pagestyle{fancy}
\fancyhf{}

\fancyhead[L]{\nouppercase{\leftmark}}
\fancyhead[R]{\materia}

\fancyfoot[C]{\thepage}

\renewcommand{\headrulewidth}{0.4pt}
\renewcommand{\footrulewidth}{0pt}

\fancypagestyle{plain}{
    \fancyhf{}
    \fancyfoot[C]{\thepage}
    \renewcommand{\headrulewidth}{0pt}
}

% =============================================================
% COMANDOS MATEMÁTICOS
% =============================================================
\newcommand{\abs}[1]{\left|#1\right|}
\newcommand{\norm}[1]{\left\lVert#1\right\rVert}
\newcommand{\vect}[1]{\mathbf{#1}}
\newcommand{\deriv}[2]{\frac{d#1}{d#2}}
\newcommand{\pderiv}[2]{\frac{\partial#1}{\partial#2}}

% =============================================================
% ENTORNOS / CAJAS ACADÉMICAS
% =============================================================
\newtcolorbox{definition}[1]{
    enhanced,
    breakable,
    title={Definición: #1},
    fonttitle=\bfseries,
    colback=blue!3,
    colframe=blue!50!black,
    boxrule=0.8pt,
    arc=2mm
}

\newtcolorbox{important}{
    enhanced,
    breakable,
    title={Importante},
    fonttitle=\bfseries,
    colback=yellow!8,
    colframe=orange!70!black,
    boxrule=0.8pt,
    arc=2mm
}

\newtcolorbox{example}[1]{
    enhanced,
    breakable,
    title={Ejemplo: #1},
    fonttitle=\bfseries,
    colback=green!4,
    colframe=green!40!black,
    boxrule=0.8pt,
    arc=2mm
}

\newtcolorbox{formula}[1]{
    enhanced,
    breakable,
    title={Fórmula / Regla: #1},
    fonttitle=\bfseries,
    colback=purple!4,
    colframe=purple!50!black,
    boxrule=0.8pt,
    arc=2mm
}

\newtcolorbox{exercise}[1]{
    enhanced,
    breakable,
    title={Ejercicio: #1},
    fonttitle=\bfseries,
    colback=gray!5,
    colframe=gray!60!black,
    boxrule=0.8pt,
    arc=2mm
}

\newtcolorbox{note}{
    enhanced,
    breakable,
    title={Nota},
    fonttitle=\bfseries,
    colback=white,
    colframe=gray!50,
    boxrule=0.6pt,
    arc=2mm
}

% =============================================================
% DOCUMENT
% =============================================================
\begin{document}

% -------------------------------------------------------------
% PORTADA
% -------------------------------------------------------------
\begin{titlepage}

\thispagestyle{empty}

\begin{center}

\vspace*{1cm}

{\large \textsc{\universidad}}

\vspace{3cm}

{\Huge \textbf{\materia}}

\vspace{1cm}

{\LARGE \linespread{1.3}\selectfont \titulo\par}

\vspace{0.8cm}

\textit{Comprender los mecanismos para resolver conflictos es aprender a construir acuerdos antes que a litigar.}

\vspace{1.2cm}

\rule{0.70\textwidth}{0.5pt}

\vspace{0.5cm}

{\large Apuntes de estudio}

\vspace{0.5cm}

\rule{0.70\textwidth}{0.5pt}

\vfill

{\large \autor}

\vspace{0.3cm}

{\large \anio}

\end{center}

\vfill

\begin{flushleft}
\textit{Este es un modesto aporte a los estudiantes de ``Derecho Procesal Civil'' no pretende sustituir el material oficial de la materia.}
\end{flushleft}

\end{titlepage}

% -------------------------------------------------------------
% DATOS DE LA CLASE (información institucional detectada en el apunte)
% -------------------------------------------------------------
\chapter*{Datos de la unidad}
\addcontentsline{toc}{chapter}{Datos de la unidad}

\begin{center}
\begin{tabularx}{0.9\textwidth}{>{\raggedright\arraybackslash}X >{\raggedright\arraybackslash}X}
\toprule
\textbf{Unidad} & Sistemas de solución de conflictos --- Métodos alternativos de resolución de conflictos \\
\textbf{Cátedra} & Prof. Dr. Rojas \\
\textbf{Docente a cargo del desarrollo} & Dra. Romina Moreno (Adjunta de Cátedra) \\
\textbf{Comisiones} & 74-41 y 74-43 \\
\bottomrule
\end{tabularx}
\end{center}

% -------------------------------------------------------------
% ÍNDICE
% -------------------------------------------------------------
\tableofcontents

% =============================================================
% CAPÍTULO 1 — TEORÍA DEL CONFLICTO: ANÁLISIS ESTÁTICO
% =============================================================
\chapter{Teoría del conflicto: análisis estático}

\section{El conflicto como materia prima del operador jurídico}

Todo abogado tiene como materia prima el conflicto. Quien se desempeña como funcionario o empleado del Poder Judicial, y en particular el juez, tiene como tarea central la resolución de conflictos. Sin embargo, en la formación universitaria la \textit{Teoría del Conflicto} no suele integrar el plan de estudios como materia obligatoria, sino que en algunas asignaturas aparece únicamente como contenido optativo.

Esta carencia formativa tiene una consecuencia práctica relevante: los operadores jurídicos tienden a tratar del mismo modo conflictos de naturaleza muy distinta ---patrimoniales, de familia, laborales, colectivos---, cuando cada tipo de conflicto requiere un mapeo específico, tanto desde su faceta estática como desde su faceta dinámica.

\begin{note}
El análisis que se desarrolla en este capítulo corresponde a la faceta \textbf{estática} de la Teoría del Conflicto. Su función es permitir que el operador identifique y mapee los elementos del conflicto antes de intervenir, de modo de contar con herramientas para gestionarlo en su faceta \textbf{dinámica} (manejo de escaladas y desescaladas).
\end{note}

También resultan relevantes para el análisis del conflicto la \textbf{Teoría de la Comunicación} ---que estudia los elementos emisor, receptor, mensaje, canal y código, y en la cual la elección del canal de comunicación puede favorecer o evitar escaladas del conflicto (por ejemplo, un WhatsApp, un email o una carta documento)--- y la \textbf{Teoría de la Información}, que permite al operador recabar y organizar los datos disponibles para generar la inteligencia necesaria en la toma de decisiones.

\section{Definición de conflicto}

En torno a la noción de conflicto existen distintas aproximaciones provenientes de la psicología, la economía y la filosofía. En el ámbito jurídico se destacan las siguientes definiciones:

\begin{definition}{Conflicto (Cossio --- Escuela Egológica del Derecho)}
Una interferencia intersubjetiva de intereses. En la conducta humana, cuando existe interferencia entre dos personas en una relación intersubjetiva cuyos intereses son incompatibles, nace el conflicto.
\end{definition}

\begin{definition}{Conflicto (Remo Entelman --- Teoría del Conflicto)}
Hay conflicto cuando dos o más personas tienen objetivos incompatibles entre sí en una relación de interdependencia.
\end{definition}

% TODO: en el apunte original no se profundiza en las razones por las cuales ambos autores son presentados como coincidentes más allá de su vinculación conceptual general.

De la definición de Entelman se derivan tres elementos necesarios para que exista conflicto:

\begin{itemize}
    \item Dos o más personas.
    \item Objetivos incompatibles entre sí.
    \item Relación de interdependencia.
\end{itemize}

\begin{important}
No es necesario que exista una norma que prohíba una conducta para que haya conflicto: dos personas pueden tener objetivos incompatibles aunque ambas realicen conductas permitidas por el derecho. Asimismo, el conflicto no debe identificarse necesariamente con la violencia: puede existir conflicto sin que se manifieste violencia alguna.
\end{important}

\section{Elementos del conflicto: análisis estático}

Identificada la definición de conflicto, corresponde analizar los elementos que lo conforman desde una perspectiva estática. Este ejercicio de mapeo, aunque el conflicto sea en esencia dinámico, permite conocer de qué tipo de conflicto se trata para poder intervenirlo y gestionarlo posteriormente.

\subsection{Los actores}

Los actores son los protagonistas del conflicto: quienes tienen una intención y una preocupación por el resultado, y que pueden generar cambios en la dinámica del conflicto en pos de obtenerlo. No se identifican necesariamente con el actor y el demandado del proceso judicial, sino con los protagonistas del conflicto concreto.

El actor del conflicto puede ser:

\begin{itemize}
    \item \textbf{Actor individual.}
    \item \textbf{Actor colectivo}: con liderazgo formal o informal; con organización o sin ella.
\end{itemize}

\begin{example}{Actores colectivos}
Un grupo de personas que ocupa un inmueble no es lo mismo que un sindicato: ambos son actores colectivos, pero cada supuesto debe analizarse específicamente según su organización y liderazgo.
\end{example}

Resulta relevante preguntarse si están presentes en la mesa de negociación todos los actores que deberían estarlo.

\begin{example}{Actores no evidentes}
En un divorcio, los actores son los cónyuges, pero la nueva pareja de una de las partes puede incidir fuertemente en el resultado ---por ejemplo, en la división del régimen de comunidad de bienes---. No necesariamente debe sentarse a esa persona en la mesa de negociación, pero sí puede resultar necesario incorporarla a la reunión entre el abogado y su cliente, en la medida en que puede condicionar el resultado.
\end{example}

\begin{definition}{Conciencia del actor}
La conciencia es el producto de un acto intelectual en el cual el actor admite encontrarse, respecto de otro actor, en una situación en la que cree ---o ambos creen--- tener objetivos incompatibles.
\end{definition}

\begin{important}
Trabajar sobre la conciencia del actor respecto de la existencia del conflicto es fundamental. Cuando un factor no tiene conciencia de estar en un conflicto, tiende a no legitimar al otro en la mesa de negociación, lo que bloquea cualquier posibilidad de gestión. Legitimar al adversario es un paso necesario para poder avanzar hacia un acuerdo.
\end{important}

\begin{example}{Falta de conciencia del conflicto}
Los trabajadores de una empresa plantean al presidente mejoras en las condiciones laborales, y amenazan con parar la planta si no son atendidos. Si el abogado le indica al presidente que no tiene obligación legal de conceder el aumento porque no se firmaron paritarias, lo está llevando a creer que no existe un conflicto, cuando en rigor sí lo hay. Trabajar sobre esa conciencia resulta fundamental para la gestión.
\end{example}

Vinculada a la conciencia se encuentra la percepción: el contenido con el que acceden al intelecto del actor las opiniones, amenazas o riesgos presentes en la realidad, que le permiten anticipar cómo podría reaccionar el otro actor frente a determinada conducta. La percepción depende en gran medida del conocimiento y la información que se tenga sobre el otro actor.

\subsection{Los objetivos}

Los objetivos son aquello que quieren los protagonistas del conflicto, lo que perciben como incompatible con el otro actor. Se distingue entre:

\begin{itemize}
    \item \textbf{Objetivos tangibles}: pueden verbalizarse y especificarse concretamente (por ejemplo: ``mi pretensión son cien mil pesos'').
    \item \textbf{Objetivos ocultos}: no se manifiestan, por razones estratégicas, de prestigio o vergüenza, o porque el actor sabe que si lo verbaliza la otra parte lo negará automáticamente según el tipo de relación o las características de su personalidad.
\end{itemize}

\subsection{Las relaciones de poder}

\begin{definition}{Poder}
En la Teoría del Conflicto, el poder se define como los recursos con los que cuenta cada uno de los actores.
\end{definition}

Resulta importante enumerar los recursos de poder del propio cliente y también los del contrario, ya que ello permite identificar con qué herramientas trabajar en la gestión de la dinámica del conflicto. Entre los recursos de poder se mencionan:

\begin{itemize}
    \item La posibilidad de generar amenaza de sanción o promesa de premio (por ejemplo, directivos de una empresa frente a sus trabajadores).
    \item El tiempo de espera para alcanzar un acuerdo (no representa lo mismo para quien no necesita cobrar de manera urgente que para quien sí lo necesita).
    \item La paciencia.
    \item La alternativa al acuerdo.
    \item La posibilidad de realizar la primera oferta.
    \item La información disponible sobre el caso y sus distintos elementos.
    \item El conocimiento que se tiene del otro actor (especialmente relevante en conflictos de familia).
\end{itemize}

\subsection{Conflictos de suma cero y de suma variable}

\begin{important}
Formar a los abogados bajo la idea de que todo lo que el otro gana se pierde uno mismo genera una mentalidad de \textbf{suma cero}, propia de los llamados \textbf{conflictos puros} (siempre hay un ganador y un perdedor). Los \textbf{conflictos impuros}, en cambio, son de \textbf{suma variable}: admiten una distribución de ganancias entre los actores.
\end{important}

Que un actor entienda que un conflicto es de suma cero está vinculado a que no cree en la posibilidad de generar soluciones que beneficien a ambas partes. Sin embargo, mediante un abordaje cooperativo y creativo pueden alcanzarse soluciones que beneficien a todos los actores involucrados; en la formación del operador de conflictos reside buena parte de esa posibilidad.

\subsection{Los terceros y las coaliciones}

Corresponde evaluar qué terceros pueden intervenir en el conflicto o incidir en su resultado, y considerar la posibilidad de establecer coaliciones o estrategias de alianza con alguno de ellos. Este es otro de los elementos que integran el análisis estático de la Teoría del Conflicto.

\section{Resolución versus solución del conflicto}

Una pregunta central en esta materia es cuál es la finalidad del derecho: la paz social, es decir, poner fin a los conflictos. Sin embargo, no siempre la sentencia pone fin al conflicto en un sentido sustancial.

\begin{important}
Debe distinguirse entre:
\begin{itemize}
    \item \textbf{Resolución del conflicto}: el acto formal ---la sentencia, o en su caso el laudo--- que pone fin al proceso.
    \item \textbf{Solución del conflicto}: la satisfacción real de los intereses de las partes, que no siempre resulta de la resolución formal.
\end{itemize}
\end{important}

\begin{example}{Conflictos de familia}
En un conflicto de familia, lo que determine el juez puede no poner fin al conflicto sustancial, sino constituir un eslabón dentro de él, dado que este tipo de conflictos incorpora elementos adicionales ---emociones, sentimientos, la cuestión familiar en sí misma--- que exceden lo resuelto formalmente.
\end{example}

De allí la importancia de que las partes, asistidas por operadores de conflictos, puedan alcanzar su propia solución, ya que en general resultará preferible a la que disponga un juez.

% -------------------------------------------------------------
% CORNELL — CAPÍTULO 1
% -------------------------------------------------------------
\section*{Repaso Cornell del capítulo}
\addcontentsline{toc}{section}{Repaso Cornell del capítulo}

\begin{table}[H]
\centering
\begin{tabularx}{\textwidth}{
    >{\raggedright\arraybackslash}p{0.30\textwidth}
    >{\raggedright\arraybackslash}X
}
\toprule
\textbf{Preguntas / palabras clave} & \textbf{Notas / respuestas} \\
\midrule

\textbf{¿Qué es el conflicto?} &
Para Cossio, una interferencia intersubjetiva de intereses. Para Entelman, la situación en que dos o más personas tienen objetivos incompatibles en una relación de interdependencia. No requiere norma prohibitiva ni implica necesariamente violencia. \\

\textbf{¿Cuáles son los elementos del análisis estático?} &
Los actores (individuales o colectivos), los objetivos (tangibles u ocultos), las relaciones de poder (recursos de cada actor) y los terceros con posibilidad de coaliciones. \\

\textbf{¿Qué es la conciencia del conflicto?} &
Acto intelectual por el cual un actor admite encontrarse en una situación de conflicto con otro. Sin conciencia no se legitima al adversario, lo que bloquea la gestión del conflicto. \\

\textbf{¿Qué diferencia a los conflictos de suma cero de los de suma variable?} &
En los de suma cero (conflictos puros) hay un ganador y un perdedor. En los de suma variable (conflictos impuros) existe una distribución de ganancias que permite soluciones cooperativas y creativas. \\

\textbf{¿En qué se diferencian resolución y solución del conflicto?} &
La resolución es el acto formal (sentencia o laudo) que pone fin al proceso. La solución es la satisfacción real de los intereses de las partes, que no siempre coincide con la resolución formal (por ejemplo, en conflictos de familia). \\

\midrule

\multicolumn{2}{p{\textwidth}}{
\textbf{Resumen:} el conflicto surge cuando dos o más personas tienen objetivos incompatibles en una relación de interdependencia, sin que sea necesaria una norma prohibitiva ni la existencia de violencia. Su análisis estático permite mapear actores, objetivos, poder y terceros antes de intervenir. Distinguir la conciencia del conflicto, el tipo de suma en juego (cero o variable) y la diferencia entre resolución formal y solución sustancial resulta indispensable para que el operador jurídico pueda gestionar el conflicto de manera adecuada.
} \\

\bottomrule
\end{tabularx}
\end{table}

% =============================================================
% CAPÍTULO 2 — SISTEMAS DE RESOLUCIÓN DE CONFLICTOS
% =============================================================
\chapter{Sistemas de resolución de conflictos}

\section{Clasificación general}

Al hablar de sistemas alternativos de resolución de conflictos, o de resolución alternativa de disputas, se distinguen dos grandes categorías.

\begin{important}
\textbf{Sistemas no adversariales} (matriz: la negociación):
\begin{itemize}
    \item \textbf{Negociación directa} (entre las partes; autocomposición del conflicto): transacción directa entre partes.
    \item \textbf{Negociación asistida} (interviene un tercero que colabora con las partes para posibilitar un acuerdo): mediación, conciliación, y otros sistemas (facilitadores, tercero neutral del sistema anglosajón, etc.).
\end{itemize}

\textbf{Sistemas adversariales} (un tercero decide sobre el conflicto):
\begin{itemize}
    \item Proceso judicial.
    \item Arbitraje: arbitraje de derecho, arbitraje de amigables componedores, pericia arbitral.
\end{itemize}
\end{important}

Puede trazarse una línea horizontal imaginaria: en un extremo se ubica la negociación directa entre las partes (nula intervención de terceros) y en el otro extremo el proceso judicial (máxima intervención de terceros para poner fin al conflicto).

\section{Comparación entre sistemas adversariales y no adversariales}

\begin{table}[H]
\centering
\caption{Diferencias entre sistemas no adversariales y adversariales}
\begin{tabularx}{\textwidth}{
    >{\raggedright\arraybackslash}p{0.25\textwidth}
    >{\raggedright\arraybackslash}X
    >{\raggedright\arraybackslash}X
}
\toprule
\textbf{Criterio} & \textbf{No adversarial} & \textbf{Adversarial} \\
\midrule
Intervención de tercero & No siempre (solo en negociación asistida) & Siempre (juez o árbitro que decide) \\
Mecanismo procedimental & Rige el principio de informalidad & Fijado por ley (proceso judicial) o por las partes (arbitraje) \\
Procedimiento probatorio & No es necesario (no hay que acreditar pretensiones) & Sí es necesario (el tercero evalúa la prueba) \\
Quién resuelve & Las propias partes se dan la solución & El juez o árbitro impone la decisión \\
Resultado & Acuerdo (con valor de sentencia si es homologado) & Sentencia o laudo arbitral \\
\bottomrule
\end{tabularx}
\end{table}

\section{Aclaración terminológica: la palabra ``alternativos''}

\begin{note}
Denominar a estos métodos como ``alternativos'' presupone la existencia de un método previo respecto del cual serían una alternativa, es decir, el proceso judicial. Sin embargo, la mediación, la conciliación y el arbitraje existen prácticamente desde el origen de la humanidad ---en tribus, clanes y sectas, y a lo largo de toda la historia del derecho, incluido el derecho romano---. Por ello, afirmar que el proceso judicial fue previo a estos métodos constituye un error histórico, y la calificación de ``alternativos'' debe tomarse con cautela.
\end{note}

% -------------------------------------------------------------
% CORNELL — CAPÍTULO 2
% -------------------------------------------------------------
\section*{Repaso Cornell del capítulo}
\addcontentsline{toc}{section}{Repaso Cornell del capítulo}

\begin{table}[H]
\centering
\begin{tabularx}{\textwidth}{
    >{\raggedright\arraybackslash}p{0.30\textwidth}
    >{\raggedright\arraybackslash}X
}
\toprule
\textbf{Preguntas / palabras clave} & \textbf{Notas / respuestas} \\
\midrule

\textbf{¿Qué caracteriza a los sistemas no adversariales?} &
Tienen como matriz la negociación (directa o asistida). No siempre interviene un tercero, rige la informalidad, no hay procedimiento probatorio y son las propias partes quienes se dan la solución. \\

\textbf{¿Qué caracteriza a los sistemas adversariales?} &
Comprenden el proceso judicial y el arbitraje. Siempre interviene un tercero que decide (juez o árbitro), el mecanismo está fijado por ley o por las partes, existe procedimiento probatorio, y el resultado es una sentencia o un laudo. \\

\textbf{¿Por qué se cuestiona el término ``alternativos''?} &
Porque supone que el proceso judicial fue el método previo, cuando en realidad la mediación, la conciliación y el arbitraje existen desde el origen de la humanidad, incluso antes que el proceso judicial tal como se lo conoce. \\

\midrule

\multicolumn{2}{p{\textwidth}}{
\textbf{Resumen:} los sistemas de resolución de conflictos se dividen en no adversariales (negociación directa o asistida, sin intervención necesaria de un tercero y con informalidad) y adversariales (proceso judicial y arbitraje, con intervención decisoria de un tercero y procedimiento probatorio). La denominación ``alternativos'' es históricamente imprecisa, dado que estos métodos preceden al proceso judicial.
} \\

\bottomrule
\end{tabularx}
\end{table}

% =============================================================
% CAPÍTULO 3 — EL MÉTODO DE HARVARD
% =============================================================
\chapter{El Método de Harvard: bases para la negociación}

En el marco de los sistemas no adversariales, cuya matriz es la negociación, existen distintas escuelas y técnicas. Entre ellas se destaca el Método de Harvard, desarrollado por los autores Fisher y Ury, de la Escuela de Derecho de la Universidad de Harvard, quienes establecieron las bases de la negociación en cuatro puntos.

\section{Separar a las personas del problema}

\begin{important}
Sea duro con el problema y suave con las personas.
\end{important}

Muchas veces, por no querer escuchar al interlocutor o por querer agredir a la persona a raíz del conflicto, se imposibilita cualquier negociación. La persona no es el problema; el problema es uno y las personas son distintas. En ocasiones, cambiar los interlocutores ---a través de los profesionales intervinientes--- mejora la comunicación. Un abogado formado en teoría del conflicto y negociación intentará mejorar la comunicación para alcanzar un acuerdo, cumpliendo así con este primer punto.

\section{Concentrarse en los intereses, no en las posiciones}

Lo primero que manifiestan los actores de un conflicto son sus posiciones, que suelen parecer inamovibles e incompatibles. Sin embargo, las posiciones no son lo que verdaderamente quieren las partes: detrás de ellas se encuentra el interés, es decir, lo que verdaderamente se persigue como objetivo, que puede resultar compatible con el de la otra parte.

\begin{example}{El conflicto de las naranjas}
Dos personas quieren la misma naranja; esa es su posición, y el conflicto parece innegociable. Sin embargo, al indagar el interés de cada parte se advierte que una quería la naranja para comerla (le interesaba el jugo y la pulpa) y la otra para hacer dulce (le interesaba la cáscara). A partir de esa diferencia puede alcanzarse un acuerdo: una cosa es la posición y otra, distinta, es el interés.
\end{example}

\section{Generar opciones: el torbellino de ideas}

Consiste en generar diversas opciones antes de tomar una decisión y elaborar una propuesta de acuerdo. Se trata de que las partes ensayen una variedad de posibilidades, abandonando la hipótesis de que necesariamente hay un ganador y un perdedor, para explorar la posibilidad de un acuerdo que beneficie a más partes. No todas las ideas generadas en este ejercicio resultan factibles; lo relevante es realizarlo para descubrir posibilidades que inicialmente no se habían imaginado.

\section{Buscar criterios objetivos}

Se trata de identificar criterios objetivos que sirvan de base al acuerdo y que contribuyan a convencer a las partes de que la solución propuesta es viable y positiva.

\begin{example}{Criterios objetivos}
Buscar casos similares resueltos en la jurisprudencia, o recurrir a uno o varios tasadores para objetivar el valor de un inmueble o de un campo en disputa.
\end{example}

\begin{note}
El criterio objetivo contribuye al convencimiento de que lo acordado es razonable. Si las partes no logran darse su propia solución, el camino que resta es el sistema adversarial.
\end{note}

% -------------------------------------------------------------
% CORNELL — CAPÍTULO 3
% -------------------------------------------------------------
\section*{Repaso Cornell del capítulo}
\addcontentsline{toc}{section}{Repaso Cornell del capítulo}

\begin{table}[H]
\centering
\begin{tabularx}{\textwidth}{
    >{\raggedright\arraybackslash}p{0.30\textwidth}
    >{\raggedright\arraybackslash}X
}
\toprule
\textbf{Preguntas / palabras clave} & \textbf{Notas / respuestas} \\
\midrule

\textbf{¿En qué consiste separar a las personas del problema?} &
En ser duro con el problema y suave con las personas, evitando que la agresión hacia el interlocutor imposibilite la negociación. Cambiar los interlocutores a través de profesionales puede mejorar la comunicación. \\

\textbf{¿Qué significa concentrarse en los intereses y no en las posiciones?} &
Ir más allá de lo que las partes manifiestan como posición (aparentemente inamovible) para identificar el interés real que la sostiene, el cual puede resultar compatible con el de la otra parte (ejemplo del conflicto de las naranjas). \\

\textbf{¿Qué es el torbellino de ideas?} &
La generación de múltiples opciones de acuerdo antes de decidir, abandonando la hipótesis de ganador-perdedor, para descubrir posibilidades de acuerdo no imaginadas inicialmente. \\

\textbf{¿Para qué sirven los criterios objetivos?} &
Para fundamentar el acuerdo en elementos externos y verificables (jurisprudencia similar, tasaciones), lo que ayuda a convencer a las partes de que la solución es razonable y viable. \\

\midrule

\multicolumn{2}{p{\textwidth}}{
\textbf{Resumen:} el Método de Harvard, desarrollado por Fisher y Ury, propone cuatro bases para la negociación: separar a las personas del problema, concentrarse en los intereses y no en las posiciones, generar opciones mediante el torbellino de ideas, y fundamentar el acuerdo en criterios objetivos. Estas bases buscan alejar la negociación de la lógica de ganador-perdedor y orientarla hacia soluciones cooperativas.
} \\

\bottomrule
\end{tabularx}
\end{table}

% =============================================================
% CAPÍTULO 4 — MEDIOS DE RESOLUCIÓN: MEDIACIÓN, CONCILIACIÓN Y ARBITRAJE
% =============================================================
\chapter{Mediación, conciliación y arbitraje}

\section{Mediación}

La mediación en el ámbito nacional está regulada por la Ley 26.589 y su decreto reglamentario N.\textsuperscript{o} 1.467/2011, normas que integran las leyes complementarias del Código Procesal y que establecen la forma en que el mediador convoca la audiencia, cómo se lo designa y cuál es el procedimiento aplicable.

\begin{important}
La mediación constituye, en el ámbito nacional, una instancia \textbf{previa y obligatoria} en los asuntos de contenido patrimonial civil y comercial, con las excepciones que la propia ley establece.
\end{important}

Una vez agotada esa instancia, se extiende un acta en la que se deja constancia de quiénes fueron las partes intervinientes, cuál era el objeto de la mediación y que no se alcanzó un acuerdo. Al no lograrse acuerdo, queda habilitada la instancia judicial para promover el juicio; el acta de cierre de la mediación previa y obligatoria constituye uno de los requisitos para promover la demanda.

\subsection{El mediador: requisitos y función}

Requisitos para ser mediador:

\begin{itemize}
    \item Ser abogado matriculado.
    \item Tener dos años de antigüedad en la matrícula.
    \item Reunir los requisitos adicionales exigidos para ser habilitado.
\end{itemize}

La habilitación y el registro de mediadores están a cargo del Ministerio de Justicia, dependencia del Poder Ejecutivo.

La función del mediador consiste en acercar a las partes y mejorar la comunicación entre ellas, de modo que sean ellas mismas quienes alcancen un acuerdo sobre sus diferencias.

\begin{important}
El mediador no debe proponer en forma directa fórmulas conciliatorias ni una posible solución al conflicto; su función es facilitar que las propias partes encuentren la solución.
\end{important}

\subsection{Principio de confidencialidad}

En el sistema de mediación rige el principio de confidencialidad: lo manifestado, dicho, o la documentación y prueba exhibida en ese ámbito, no puede ser utilizado en un proceso judicial posterior. Ello permite que las partes negocien con la libertad suficiente para abordar las distintas aristas del conflicto. El mediador no puede ser convocado como testigo en un proceso judicial posterior para relatar lo manifestado por las partes durante la mediación.

\subsection{Mediación oficial y mediación privada}

\begin{itemize}
    \item \textbf{Mediación oficial}: se solicita ante la cámara del fuero correspondiente el sorteo de un mediador.
    \item \textbf{Mediación privada}: se elige un mediador de confianza adecuado para el caso concreto, quien convoca a la otra parte mediante carta documento. Notificada la parte requerida, esta puede elegir entre cinco mediadores ofrecidos (incluido el ya elegido por el requirente); si no elige ninguno, queda designado el mediador propuesto por el requirente, y si elige uno del listado, queda designado ese.
\end{itemize}

% TODO: el apunte no aclara con precisión si esta posibilidad de elección entre cinco mediadores rige únicamente para la mediación privada o también admite matices en la mediación oficial; se conserva tal como fue expuesto en la clase.

\begin{note}
Esta posibilidad de elección entre varios mediadores no está prevista en la conciliación laboral previa.
\end{note}

\section{Conciliación}

\subsection{Conciliación intra-procesal}

El artículo 360 del Código Procesal Civil y Comercial, modificado por la Ley de Mediación 26.589, prevé que los jueces, cuando adviertan que el conflicto podría resolverse mediante alguno de los métodos alternativos, pueden derivar el caso a mediación o solicitar la reapertura de la mediación, dentro de un plazo de treinta días.

\begin{formula}{Art. 360 CPCC --- Audiencia preliminar}
El juez invita a las partes a alcanzar una conciliación como una de las etapas de la audiencia preliminar, la cual tiene lugar luego de transcurrida la etapa introductoria del proceso (demanda, contestación de demanda y prueba ofrecida).
\end{formula}

\begin{formula}{Art. 36 inc. 2.\textsuperscript{o} CPCC}
En cualquier momento del proceso, antes de dictada la sentencia (etapa introductoria, probatoria o conclusional), el juez puede convocar a las partes para intentar una conciliación.
\end{formula}

A diferencia del mediador, el juez, en virtud de su propia función, sí tiene la posibilidad de proponer fórmulas conciliatorias sobre la base de criterios objetivos, de la causa y de las cuestiones de prueba que se hayan suscitado, sin que ello implique un prejuzgamiento, ya que el mecanismo está previsto precisamente para evitarlo.

\begin{table}[H]
\centering
\caption{Diferencias entre el juez conciliador y el mediador}
\begin{tabularx}{\textwidth}{
    >{\raggedright\arraybackslash}p{0.30\textwidth}
    >{\raggedright\arraybackslash}X
    >{\raggedright\arraybackslash}X
}
\toprule
\textbf{Aspecto} & \textbf{Juez como conciliador} & \textbf{Mediador} \\
\midrule
Confidencialidad & No rige & Sí rige \\
Informalidad & No rige; deben cumplirse las formas del proceso & Sí rige \\
Dependencia institucional & Poder Judicial & Ministerio de Justicia (Poder Ejecutivo) \\
Proponer fórmulas conciliatorias & Sí puede & No puede \\
Nombramiento / requisitos & Es el juez de la causa & Abogado matriculado con dos años de antigüedad, habilitado por el Ministerio de Justicia \\
\bottomrule
\end{tabularx}
\end{table}

\subsection{La conciliación como modo anormal de terminación del proceso}

El modo normal de terminación del proceso es la sentencia. Uno de los modos anormales previstos por el Código Procesal es la conciliación: cuando en la audiencia preliminar, o en otro momento del proceso, se alcanza un acuerdo, este tiene la misma autoridad que una sentencia.

\begin{formula}{Art. 309 CPCC}
Los acuerdos conciliatorios celebrados por las partes ante el juez y homologados por este tendrán autoridad de cosa juzgada. Si no se cumplen, se ejecutan por la vía de ejecución de sentencia.
\end{formula}

De igual modo, el acuerdo alcanzado en el marco de una mediación ---volcado en el acta correspondiente, con sus deberes, obligaciones y consecuencias ante el incumplimiento--- tiene para las partes valor de sentencia. En caso de incumplimiento, se ejecuta por la vía de ejecución de sentencias (arts. 499 y siguientes CPCC; en particular, para los acuerdos de mediación, art. 500 CPCC).

\subsection{SECLO: conciliación laboral obligatoria}

Para los conflictos con raíz en el derecho del trabajo ---trabajadores y empleadores, diferencias salariales, sindicatos y demás cuestiones propias del derecho laboral--- corresponde transitar la conciliación laboral obligatoria, prevista y regulada a través del SECLO (Servicio de Conciliación Laboral Obligatoria), que depende del Ministerio de Trabajo.

\begin{formula}{Ley 24.635}
Crea el SECLO y modifica la Ley de Procedimiento Laboral N.\textsuperscript{o} 18.345. El acta de cierre de la conciliación previa y obligatoria constituye un requisito para promover el juicio laboral.
\end{formula}

A diferencia del mediador, el conciliador laboral sí puede proponer fórmulas conciliatorias, de modo análogo al juez. Además, debe controlar el acuerdo celebrado ante sí a los fines de su homologación por el Ministerio de Trabajo, dado que se encuentran comprometidos el orden público y la protección de los derechos laborales.

\begin{note}
La materia civil y comercial (y algunas cuestiones de índole federal) transita la mediación; la materia laboral transita la conciliación laboral obligatoria. Ambas son instancias previas y obligatorias a la promoción del juicio.
\end{note}

\section{Arbitraje}

El arbitraje se ubica dentro de los sistemas adversariales de resolución de conflictos: un tercero, el árbitro, interviene y decide sobre el conflicto.

\subsection{Materia disponible para el arbitraje}

Las partes solo pueden someter al arbitraje una materia que resulte transable y que no se encuentre afectada por el orden público. No puede someterse a arbitraje, por ejemplo, un conflicto de alimentos o de divorcio.

\subsection{La cláusula compromisoria}

El origen del arbitraje es contractual: las partes acuerdan que, en caso de conflicto derivado del contrato, la diferencia será sometida a arbitraje. Esta previsión se incorpora mediante la \textbf{cláusula compromisoria}, a través de la cual las partes se comprometen a que su conflicto sea dirimido ante un arbitraje.

\subsection{Tipos de arbitraje}

\begin{itemize}
    \item \textbf{Arbitraje ad hoc}: se crea específicamente para resolver un conflicto concreto; las partes pueden designar árbitros (por ejemplo, dos árbitros que a su vez eligen a una tercera persona, conformando un tribunal de tres árbitros).
    \item \textbf{Arbitraje institucional}: se somete el conflicto a un tribunal de arbitraje institucional ya existente (por ejemplo, el Tribunal de Arbitraje de la Bolsa de Comercio de Buenos Aires, el Tribunal de Arbitraje de la Bolsa de Cereales, o el Tribunal de Arbitraje del Colegio de Abogados de San Isidro), que cuenta con reglamento propio y árbitros permanentes designados por concurso.
\end{itemize}

\begin{formula}{Regulación del arbitraje}
El arbitraje se encuentra regulado en el Código Procesal Civil y Comercial a partir del artículo 736 y siguientes.
\end{formula}

\subsection{Modalidades de arbitraje}

\begin{definition}{Arbitraje de derecho}
El árbitro resuelve en los mismos términos en que podría resolver un juez, pero ante una jurisdicción privada, aplicando la norma. Es, en general, más rápido y económico en cuanto a costas y gastos, y constituye el método de resolución de conflictos más utilizado a nivel internacional. Su decisión se denomina \textbf{laudo}, equiparable a una sentencia: ante incumplimiento, se ejecuta judicialmente por la vía de ejecución de sentencia (art. 499 y siguientes CPCC). Los árbitros no tienen la posibilidad de ejecutar el laudo por sí mismos, para lo cual debe acudirse a la justicia.
\end{definition}

\begin{definition}{Arbitraje de amigables componedores}
No se resuelve con la misma rigidez del derecho, sino conforme a la equidad. Los árbitros tienen en cuenta la norma y el derecho, pero aplican criterios de equidad. Según Barrios de Angelis, en este tipo de arbitraje se produce la ``dulcificación de la letra de la ley'': una forma de morigerar cuestiones fijadas taxativamente en la normativa, para resolver conforme a la equidad del caso concreto.
\end{definition}

\begin{definition}{Pericia arbitral (art. 773 CPCC)}
Prevista en el artículo 773 del CPCC, que remite al artículo 516. Procede cuando existe dificultad en las operaciones a realizar por resultar las liquidaciones muy complejas, para lo cual se requiere la intervención de especialistas en la materia. Se designa un experto que emite un dictamen a los fines de acercar o dirimir las diferencias.
\end{definition}

% TODO: el apunte original plantea como interrogante abierto, atribuido al Dr. Rojas, si la pericia arbitral puede considerarse otro modo de conclusión del proceso; la clase no cierra esta cuestión y remite a un trabajo específico del autor.

% -------------------------------------------------------------
% CORNELL — CAPÍTULO 4
% -------------------------------------------------------------
\section*{Repaso Cornell del capítulo}
\addcontentsline{toc}{section}{Repaso Cornell del capítulo}

\begin{table}[H]
\centering
\begin{tabularx}{\textwidth}{
    >{\raggedright\arraybackslash}p{0.30\textwidth}
    >{\raggedright\arraybackslash}X
}
\toprule
\textbf{Preguntas / palabras clave} & \textbf{Notas / respuestas} \\
\midrule

\textbf{¿Qué régimen tiene la mediación y quién la designa?} &
Es instancia previa y obligatoria en materia civil y comercial nacional (Ley 26.589 y Decreto 1.467/2011). La designa el Ministerio de Justicia. El mediador acerca a las partes pero no propone fórmulas; rige la confidencialidad, y el acuerdo tiene valor de sentencia (art. 500 CPCC). Puede ser oficial o privada. \\

\textbf{¿Cómo interviene el juez como conciliador?} &
A través de los arts. 360 (audiencia preliminar) y 36 inc. 2.\textsuperscript{o} CPCC, en cualquier momento antes de la sentencia. A diferencia del mediador, el juez sí puede proponer fórmulas conciliatorias; no rige confidencialidad ni informalidad. El acuerdo logrado tiene autoridad de cosa juzgada (art. 309 CPCC). \\

\textbf{¿Qué es el SECLO y cuándo se aplica?} &
Es el Servicio de Conciliación Laboral Obligatoria (Ley 24.635), instancia previa y obligatoria en materia laboral, dependiente del Ministerio de Trabajo. El conciliador puede proponer fórmulas, y el acuerdo debe ser homologado por dicho Ministerio. \\

\textbf{¿Cuándo procede el arbitraje y cómo nace?} &
Procede sobre materia disponible, no afectada por el orden público (no en alimentos o divorcio). Nace de una cláusula compromisoria incorporada al contrato. Puede ser ad hoc o institucional. \\

\textbf{¿Cuáles son las modalidades de arbitraje?} &
Arbitraje de derecho (resuelve conforme a la norma; la decisión es el laudo), arbitraje de amigables componedores (resuelve según la equidad) y pericia arbitral (art. 773 CPCC, para liquidaciones complejas). \\

\midrule

\multicolumn{2}{p{\textwidth}}{
\textbf{Resumen:} la mediación (Ley 26.589) y la conciliación laboral (SECLO, Ley 24.635) son instancias previas y obligatorias, la primera en materia civil y comercial y la segunda en materia laboral; en ambas puede proponerse un acuerdo con valor de sentencia. La conciliación judicial (arts. 360 y 36 inc. 2.\textsuperscript{o} CPCC) permite al juez proponer fórmulas conciliatorias, con acuerdo homologado bajo autoridad de cosa juzgada (art. 309 CPCC). El arbitraje, por su parte, opera sobre materia disponible mediante cláusula compromisoria, y puede ser de derecho, de amigables componedores o pericia arbitral.
} \\

\bottomrule
\end{tabularx}
\end{table}

% =============================================================
% CAPÍTULO 5 — REFERENCIAS BIBLIOGRÁFICAS Y NORMATIVAS
% =============================================================
\chapter{Referencias bibliográficas y normativas}

\section{Bibliografía recomendada}

\begin{itemize}
    \item Entelman, Remo: \textit{Teoría del Conflicto} (análisis estático y dinámico completo del conflicto).
    \item Calvo Soler, Raúl: \textit{Mapeo de Conflictos} (proyecciones de casos concretos con análisis estático y dinámico).
    \item Rojas: trabajos sobre arbitraje, disponibles en \url{www.rojasomar.com.ar}, especialmente ``La pericia arbitral como otro modo de conclusión del proceso''.
    \item Fisher, R. y Ury, W.: Método de Harvard sobre negociación.
\end{itemize}

\section{Normativa citada}

\begin{longtable}{
    >{\raggedright\arraybackslash}p{0.28\textwidth}
    >{\raggedright\arraybackslash}p{0.62\textwidth}
}
\toprule
\textbf{Norma} & \textbf{Contenido / función} \\
\midrule
\endfirsthead
\toprule
\textbf{Norma} & \textbf{Contenido / función} \\
\midrule
\endhead
\bottomrule
\endfoot
\bottomrule
\endlastfoot

Ley 26.589 & Mediación prejudicial obligatoria; rige para materia civil y comercial en el ámbito nacional. \\
Decreto 1.467/2011 & Decreto reglamentario de la Ley 26.589. \\
Ley 24.635 & Crea el SECLO. Modifica la Ley de Procedimiento Laboral N.\textsuperscript{o} 18.345. Conciliación laboral obligatoria. \\
Art. 36 inc. 2.\textsuperscript{o} CPCC & Facultad del juez de convocar a las partes para intentar una conciliación en cualquier momento del proceso. \\
Art. 360 CPCC & Audiencia preliminar; invitación a conciliación; derivación a mediación (modificado por la Ley 26.589). \\
Art. 309 CPCC & Acuerdos conciliatorios con autoridad de cosa juzgada (modo anormal de terminación del proceso). \\
Arts. 499 y ss. CPCC & Ejecución de sentencia; base para ejecutar laudos arbitrales y acuerdos de mediación. \\
Art. 500 CPCC & Específicamente, la vía de ejecución del acuerdo de mediación. \\
Arts. 736 y ss. CPCC & Regulación del arbitraje en el Código Procesal Civil y Comercial. \\
Art. 773 CPCC & Pericia arbitral (remite al art. 516). \\
Art. 516 CPCC & Liquidaciones complejas (referenciado por el art. 773 sobre pericia arbitral). \\

\end{longtable}

\end{document}
